\section{Introduction}
A free surface interrupts bulk coordination and carries an excess energy that depends on crystallographic orientation, chemistry, and atomic arrangement. In disordered alloys, surface and subsurface sites can have local compositions different from the bulk, so nominally identical samples may yield different energies. Slab dimensions and bulk-reference consistency must therefore be checked before comparing facets \cite{sun2013slabs,yalcin2024surfflow}.

This study uses LAMMPS molecular dynamics and a modified embedded-atom method (MEAM) potential to compare FCC-based FeNiAl, Fe$_{20}$Ni$_{50}$Al$_{30}$, and Fe$_{50}$Ni$_{30}$Al$_{20}$ at 300, 500, and 700~K. The group dataset covers (100), (110), and (111) surfaces, with three chemical configurations per condition. The structure-generation procedure imposes FCC; it does not establish FCC as the stable phase at every composition and temperature.
